\section{Introduction}
\label{sec:intro}

WireGuard~\cite{donenfeld2017wireguard} is a VPN protocol whose handshake
instantiates the Noise \texttt{IKpsk2} pattern~\cite{perrin2018noise}. It sits
in the mainline Linux kernel and carries production traffic, so questions
about its security properties, and about what migrating it to post-quantum
primitives would cost, are engineering questions with deployed consequences.

The protocol has been analysed before: symbolically by Donenfeld and
Milner~\cite{donenfeld2018formal}, and in the computational model by Lipp,
Blanchet and Bhargavan~\cite{lipp2019cryptoverif}. This work is an independent
reconstruction from the specification and claims no new vulnerability. Its
contribution lies elsewhere, in two places that prior analyses did not
occupy. The first is methodological: the handshake is transcribed \emph{once}
and handed to two provers, so that a disagreement between them is evidence
about the provers rather than about two people's modelling habits. The second
is that the analysis does not stop at the cryptography. Post-quantum
parameter sizes are followed through to datagram sizes and from there to path
MTUs, which is where the migration cost actually becomes visible.

\begin{figure*}[t]
\centering
\resizebox{\textwidth}{!}{\begin{tikzpicture}[x=1cm,y=1cm,font=\footnotesize,
    bx/.style={draw=black!55,rounded corners=2pt,align=center,inner sep=4pt,
               text width=3.2cm,minimum height=1.0cm},
    hi/.style={bx,draw=okgreen,very thick},
    lnk/.style={->,thick,black!45}]

  \node[bx] (spec) at (1.7,0)
    {\textbf{WireGuard specification}\\{\scriptsize wire format $+$ Noise \texttt{IKpsk2}}};
  \node[hi] (core) at (5.6,0)
    {\textbf{one shared core}\\{\scriptsize \texttt{wireguard\_core.pvl}}};
  \node[bx] (pv) at (9.5,1.05)
    {\textbf{ProVerif models}\\{\scriptsize one file per property}};
  \node[bx] (tam) at (9.5,-1.05)
    {\textbf{Tamarin theory}\\{\scriptsize transcribed term for term}};
  \node[bx] (run) at (13.4,0)
    {\textbf{\texttt{runner.py}}\\{\scriptsize wall-clock $+$ memory budget}};
  \node[bx] (raw) at (17.3,0)
    {\textbf{\texttt{results/raw/}}\\{\scriptsize verbatim prover output}};

  \node[hi] (json) at (17.3,-3.1)
    {\textbf{\texttt{results.json}}\\{\scriptsize machine-readable verdicts}};
  \node[bx] (exp) at (13.4,-3.1)
    {\textbf{\texttt{expectations.yaml}}\\{\scriptsize regression gate}};
  \node[bx] (gen) at (9.5,-3.1)
    {\textbf{\texttt{render\_tables.py}}\\{\scriptsize tables, figures, \texttt{facts.tex}}};
  \node[bx] (tex) at (5.6,-3.1)
    {\textbf{\texttt{report/main.tex}}\\{\scriptsize prose only; no typed numbers}};
  \node[hi] (pdf) at (1.7,-3.1)
    {\textbf{\texttt{docs/report.pdf}}\\{\scriptsize built by \texttt{make report}}};

  \draw[lnk] (spec) -- (core);
  \draw[lnk] (core.east) -- (pv.west);
  \draw[lnk] (core.east) -- (tam.west);
  \draw[lnk] (pv.east) -- (run.west);
  \draw[lnk] (tam.east) -- (run.west);
  \draw[lnk] (run) -- (raw);
  \draw[lnk] (raw) -- node[left,black!55,font=\scriptsize,align=right]
    {\texttt{collect\_}\\\texttt{results.py}} (json);
  \draw[lnk] (json) -- (exp);
  \draw[lnk] (exp) -- (gen);
  \draw[lnk] (gen) -- (tex);
  \draw[lnk] (tex) -- (pdf);
\end{tikzpicture}
}
\caption{How one protocol description reaches two provers and one report.
The two shaded stages are the load-bearing ones: a single shared core that
both prover front-ends include, and a single machine-readable result file that
every table, figure and quoted number in this paper is generated from. No
number in this document is typed by hand.}
\label{fig:pipeline}
\end{figure*}

\subsection{Contributions}

\para{A controlled experiment isolating the pre-shared key.} The \texttt{psk}
of \texttt{IKpsk2} is widely described as providing post-quantum protection.
Section~\ref{sec:results} makes that claim testable with a pair of models
identical in every respect except whether the pre-shared key is secret, run
against an attacker that gains a discrete-logarithm oracle after the session
completes. Transport-key secrecy is proved in one and refuted in the other.
Because the models differ in exactly one variable, so does the explanation.

\para{Locating key-compromise impersonation resistance.} With the responder's
own static key in the attacker's hands, we show that the responder \emph{does}
commit transport keys for a session it attributes to an honest initiator, and
\emph{never} accepts traffic on that session. Forging message~1 needs only
Diffie-Hellman terms computable from the leaked scalar; deriving the transport
keys additionally needs the static-ephemeral term \(se=\mathrm{DH}(e_r,S_i)\),
which does not follow. The property therefore lives in the interval between
committing key material and accepting traffic---a distinction a single
monolithic ``session established'' event would hide entirely.

\para{A three-dimensional account of post-quantum migration cost.}
Section~\ref{sec:pq} follows a chain of consequences: KEM parameter sizes
determine datagram sizes, datagram sizes meet path MTUs, and a datagram that
exceeds the IPv6 minimum MTU has no guaranteed path. An in-band ML-KEM-768
initiation is \wgKemInit\,B against a usable payload of \wgIpvsixPayload\,B.
The cost of migration is not only computational: it appears simultaneously in
packet size, in path reachability, and in the strength of the binding
assumption the KEM must satisfy.

\para{Reported failure, and numbers that cannot drift.} Of \wgPvProps{}
ProVerif properties, \wgPvNonterm{} did not terminate within the budget this
project set itself, and two of them were decided by Tamarin in seconds. Both facts are in the result
table alongside the proofs. So is the mechanism that keeps them honest: the
tables, the figures and every quantity quoted in this text are generated from
\texttt{results/results.json} at build time (Figure~\ref{fig:pipeline}), so a
sentence in this paper cannot survive a run that contradicts it.

\subsection{Scope}

This paper concerns the handshake. The data plane, the cookie-based
rate-limiting mechanism, and endpoint roaming are network-layer behaviours
with no counterpart in an abstract Noise model;
Section~\ref{sec:modelling} states for each wire-format field whether it is
represented in the models and, where it is not, why. Reproduction
instructions and the resource profile of every result accompany the artifact.
